\documentclass[11pt,letterpaper]{article}
\usepackage[utf8]{inputenc}
\usepackage[T1]{fontenc}
\usepackage{amsmath,amsfonts,amssymb}
\usepackage{graphicx}
\usepackage{geometry}
\usepackage{subcaption}
\usepackage{caption}
\usepackage{hyperref}

\geometry{margin=1in}
\setlength{\parskip}{4pt plus 1pt}
\setlength{\parindent}{0pt}

\begin{document}

\noindent\textbf{1.5 Mechano-electrochemical simulation and digital
twin development (RIT)}

The RIT team will lead theoretical efforts, including AI-enabled
atomistic calculations of key materials involved in the cell design
and an interpretable AI-driven electro--chemo--mechanical digital twin
that quantitatively determines design parameters.

\textbf{Subtask I: Material-parameter calculation.}
AI-enabled DFT/MD calculations will estimate adhesion energies for
LVPF/PTFE, graphite/PTFE, and electrode/GPE interfaces, which define
cohesive parameters for debonding and contact-retention modeling.
MD/coarse-grained simulations will compute GPE modulus, ionic
conductivity, and relaxation time as functions of fluorinated-acrylate
and liquid-electrolyte fractions.  Continuum composite calculations
will estimate package shrink stress and barrier-layer stiffness as
functions of shrink strain, BN loading, coating modulus, and coating
thickness.

\textbf{Subtask II: Digital twin development (Fig.~\ref{fig:dt}a).}
A coupled chemo-mechanical model integrates Subtask~I parameters and
experiment with the cell-level wrapped-pouch geometry.
Polymerization-induced eigenstrains in the GPE
($\varepsilon_{\rm GPE}$) and the conformal package
($\varepsilon_{\rm pkg}$), together with lithiation-driven electrode
breathing and any external stack pressure $\sigma_{\rm stack}$,
generate the through-thickness contact pressure $\langle p_c\rangle$.
A pressure-dependent contact resistance
$R_c(p_c) = R_0\,e^{-\alpha p_c}$ closes the model and yields the
central output, the contact-stability metric
$\Pi_c = \langle p_c\rangle/p_{\min}$, where $p_{\min}$ is the stress
required to maintain $R_c \le 10$--$20$~m$\Omega\cdot$cm$^2$.  Design
maps over $(\varepsilon_{\rm pkg}, \varepsilon_{\rm GPE}, E_{\rm sep},
\sigma_{\rm stack})$, with secondary axes for PTFE fraction (3--8\%),
CNT loading (1--5 wt\%), collector porosity (20--60\%), and electrode
thickness, identify operating windows where $\Pi_c \ge 1$, ionic
conductivity $\ge 1$~mS\,cm$^{-1}$, and inactive mass $\le 15\%$ are
simultaneously met without rigid compression.

\textbf{Preliminary result.}
Figure~\ref{fig:dt}b maps $\Pi_c$ at the pressure-free condition
($\sigma_{\rm stack}=0$): package shrink $\varepsilon_{\rm pkg} \gtrsim
0.07$ alone clears $\Pi_c = 1$ without external compression.  Three
actionable findings emerge.  (i)~$\varepsilon_{\rm pkg}$ is the
dominant linear lever, with $\Pi_c \propto \varepsilon_{\rm pkg}$ at
fixed $\varepsilon_{\rm GPE}$.  (ii)~$\varepsilon_{\rm GPE}$ must be
co-designed with $\varepsilon_{\rm pkg}$: high $\varepsilon_{\rm GPE}$
at low $\varepsilon_{\rm pkg}$ locally unloads the separator--cathode
interface and risks debonding.  (iii)~Stiffening the GPE
(Fig.~\ref{fig:dt}c) provides a useful but diminishing-returns boost
that saturates near $E_{\rm sep} \approx 25$~MPa.  Together these
identify an engineering window---$\varepsilon_{\rm pkg} \in [0.07,
0.10]$, $\varepsilon_{\rm GPE} \le 0.05$, $E_{\rm sep} \ge 10$~MPa---
that meets contact stability without rigid compression.

\begin{figure}[h]
  \centering
  \begin{subfigure}[b]{0.32\textwidth}
    \centering
    \includegraphics[width=\textwidth]{figures/geometry_schematic.pdf}
    \caption{}
    \label{fig:dt-a}
  \end{subfigure}\hfill
  \begin{subfigure}[b]{0.32\textwidth}
    \centering
    \includegraphics[width=\textwidth]{figures/prop_heatmap_sigma0.pdf}
    \caption{}
    \label{fig:dt-b}
  \end{subfigure}\hfill
  \begin{subfigure}[b]{0.32\textwidth}
    \centering
    \includegraphics[width=\textwidth]{figures/prop_E_sep_sigma0.pdf}
    \caption{}
    \label{fig:dt-c}
  \end{subfigure}
  \caption{(a)~Wrapped-pouch cross-section: polymer envelope (hatched)
  wraps the active stack and supplies through-thickness compression
  via cure shrinkage.  (b)~$\Pi_c$ over
  $(\varepsilon_{\rm pkg},\,\varepsilon_{\rm GPE})$ at
  $\sigma_{\rm stack}=0$, $E_{\rm sep}=10$~MPa; white contour marks
  $\Pi_c=1$.  (c)~$\Pi_c$ vs.\ GPE modulus at
  $\varepsilon_{\rm GPE}=0.02$; diminishing returns above $E_{\rm sep}
  \approx 25$~MPa.}
  \label{fig:dt}
\end{figure}

\end{document}
